%% Stand 18. Juli 2026
Dieses Kapitel ist wie folgt strukturiert:
\begin{itemize}[topsep=2pt,parsep=1pt,itemsep=0.5pt]
\item[$\circ$] Im Kapitel \ref{KapitelTheorieKernidee} wird das
Grundprinzip des Kriging als einer Methode der Vorhersage von
Zufallsvariablen im Rahmen der räumlichen Statistik erläutert.
\item[$\circ$] Das Kapitel \ref{KapitelTheorieGrundlagen} stellt 
die notwendigen Begriffe und Zusammenhänge für das weitere Verständnis bereit.
\item[$\circ$] Das Variogramm ist das zentrale Werkzeug der Modellierung
im Zuge des Kriging und wird in Kapitel \ref{KapitelTheorieVariogramm}
erläutert.
\item[$\circ$] Das Variogramm seinerseits dient als
Schätzfunktion für Kovarianz-Funktionen, welche die
Verbundenheit der Zufallsvariablen charakterisieren. Die 
gebräuchlichsten Koavrianz-Funktionen werden in
Kapitel \ref{KapitelTheorieKovarianzFunktionen} behandelt, bevor 
dann  
\item[$\circ$] in Kapitel \ref{KapitelTheorieKriging} 
die Kriging Gleichungen motiviert und aufgeführt werden.  
\item[$\circ$] Abschließend wird in Kapitel \ref{KapitelTheorieAnisotropie} 
das praxisrelevante Thema der Anisotropie hinsichtlich 
Phänomenologie, Detektion und Handhabung besprochen.  
\end{itemize}

\subsection{Grundidee des Kriging}
\label{KapitelTheorieKernidee}
Das von dem Geografen Waldo Tobler formulierte erste 
Gesetz\footnote{Ob es sich hierbei im strengen Sinne um ein Gesetz handelt
oder um eine in der Geophysik bewährte Heuristik wird 
z.\,B.\,in \cite{Miller2004} thematisiert. Die Antwort auf diese Frage ist
allerdings für die weiteren Überlegungen unerheblich.} der Geografie lautet 
(vgl.\,\cite[S.\,236]{Tobler1970}):
\begin{quote}
\glqq \textit{Everything is related to everything else, but near things are more related than distant things.}\grqq
\end{quote}

Diese Aussage wurde von Matheron (vgl.\,\cite{Matheron1971}) folgendermaßen in die Sprache
der Mathematik übersetzt:
\begin{itemize}[topsep=2pt,parsep=1pt,itemsep=0.5pt]
\item[a)] unter \textit{Dingen} (als Übersetzung für \glqq everything\grqq )
werden numerische Größen in Räumen verstanden, die mit einer Metrik 
versehen sind, beispielsweise also Messwerte an Orten auf 
der Erdoberfläche;
\item[b)] diese Größen werden als Realisierungen 
von Zufallsvariablen interpretiert;
\item[c)] die \textit{Verbundenheit} (als Übersetzung für 
\glqq related to\grqq ) dieser Zufallsvariablen wird durch 
Beziehungen ihrer Erwartungswerte und Kovarianzen ausgedrückt.
\end{itemize}

Kriging ist eine Methode zur Vorhersage von Werten an Orten, 
an denen \textit{keine} Beobachtungen vorliegen.  
Um den Vorhersageprozess mathematisch handhabbar zu machen, 
beschränkt Kriging die Klasse der Prädiktoren auf 
\textit{lineare Funktionen} der beobachteten Werte,
wodurch das Anpassungsproblem auf die Bestimmung einer endlichen
Anzahl von Koeffizienten reduziert wird. \\

Das dahinterstehende Optimierungsproblem wird durch die
Festlegung folgender Nebenbedingungen definiert:

\begin{itemize}[topsep=2pt,parsep=1pt,itemsep=0.5pt]
\item[1) ] der mittlere quadratische Vorhersagefehler sollte minimiert werden, das heißt, der Prädiktor sollte \textit{präzise} sein; 
\item[2) ] der Prädiktor sollte im Durchschnitt dem wahren Wert
entsprechen, das heißt, er sollte \textit{unverzerrt} sein.
\end{itemize}

Diese Anforderungen führen schließlich auf ein lineares Gleichungssystem mit 
einer \textit{eindeutigen} Lösung. 
Der resultierende Prädiktor wird daher als \textbf{B}est 
\textbf{L}inear \textbf{U}nbiased \textbf{P}redictor oder auf
deutsch \glqq bester linearer unverzerrter Prädiktor\grqq\ (kurz BLUP)
bezeichnet.\\
Eine Voraussetzung für die Bestimmung der optimalen Gewichte des
Gleichungssystems besteht darin, die 
räumliche Abhängigkeit der Beobachtungen quantitativ zu beschreiben. 
Genau hierzu dienen die im folgenden Kapitel eingeführten 
Kovarianz-Funktionen und die Semivarianz.\\

Das skizzierte Grundprinzip des Kriging wird in Kapitel 
\ref{KapitelTheorieKriging} näher ausgearbeitet.


